\chapter{Computational Study}
\label{sec:comp_study}

To demonstrate the practical applicability of the proposed \algoNotation framework, we apply it to airline schedule planning---a canonical instance of service network design with passenger choice. In this setting, the network nodes correspond to airports, the service resource types correspond to aircraft fleet types, the service arcs correspond to scheduled flights, and the spatial segments correspond to origin-destination airport pairs. This application domain is particularly well-suited for evaluating our framework because it exhibits all the computational challenges discussed in the preceding sections: a large number of feasible itineraries, strong coupling between markets through shared network capacity, and the need for precise temporal resolution.

We evaluate the performance of the \algoNotation framework using a real-world instance from Frontier Airlines (F9), and compare against a baseline of solving the integrated model directly using a state-of-the-art solver (referred to as the ``Raw Model'').

\section{Experimental Setup}
The Frontier Airlines (F9) case study represents an ultra-low-cost carrier (ULCC) network characteristic of highly competitive airline markets in the United States. The problem instance is constructed with a 24-hour planning horizon.

The network encompasses $60$ airports and $259$ directed flight segments. The airline operates a homogeneous Airbus A320-family fleet, modeled with $4$ specific fleet sub-types to account for varying seating capacities and operating costs. The seat capacities primarily range between $186$ and $228$ seats. We utilize a base time discretization of $15$ minutes ($96$ intervals per day). This resolution is chosen to align with the industry-standard on-time performance threshold: in commercial aviation, an arrival within $15$ minutes of the scheduled time is considered on-time by regulatory authorities. This discretization step therefore represents the finest temporal granularity at which schedule deviations are operationally meaningful.

The passenger demand side includes $382$ origin-destination (OD) markets. We consider routing options with up to $2$ connections and a maximum connection window of $3$ hours. A total of $475,776$ feasible itineraries are generated by this process. Coupling nearly half a million itinerary variables with discrete fleet assignment constraints results in an extremely dense mathematical program, pushing the limits of conventional computational resources.

\section{Computational Performance}
We compare the convergence behavior and tractability of the \algoNotation algorithm against the Raw Model baseline over a 12-hour computational limit. The results underscore the critical necessity of path-signature-based aggregation and dynamic variable refinement. The performance metrics are summarized in Table\ref{tab:comp_results}.

\begin{table}[h]
\centering
\caption{Computational Performance: Raw Model vs. \algoNotation}
\label{tab:comp_results}
\begin{tabular}{lrrrr}
\hline
Method & Objective (UB) & Lower Bound (LB) & Optimality Gap & Runtime (h) \\ \hline
Raw Model (Baseline) & $-176,824.17$ & $-185,214.71$ & $4.75\%$ & $12.00^*$ \\
\textbf{\algoNotation (Proposed)} & $\mathbf{-176,843.71}$ & $\mathbf{-176,843.71}$ & \textbf{$0.00\%$} & \textbf{10.92} \\ \hline
\multicolumn{4}{l}{\small $^*$ Terminated upon reaching the 12-hour (43,200 seconds) time limit.}
% \multicolumn{4}{l}{$\text{Gap}=\frac{|\text{UB}-\text{LB}|}{\max\{|\text{UB}|,|\text{LB}|\}}$}
\end{tabular}
\end{table}

When the complete integrated model with 475,776 itinerary variables is passed directly to the Gurobi solver, the root relaxation and integer bounding processes suffer from massive fractional degeneracy. After a continuous 12-hour ($43,200$ seconds) execution, the Raw Model only manages to find an incumbent objective of $-176,824.17$, leaving a heavily stalled optimality gap of $4.75\%$. In further analysis, the branch-and-bound tree exploration keeps expanding without any improvement in the incumbent solution or tightening of the lower bound for hours, indicating that the solver is effectively stuck in a computational quagmire due to the sheer size and complexity of the original formulation. The unexplored node counts (the number of nodes that have been generated but not yet explored) of branch-and-bound trees of models are plotted in Figure\ref{fig:comp_tree}, which clearly shows that the Raw Model's branch-and-bound tree keeps expanding, meanwhile the incumbent solution and the lower bound remain nearly unchanged. 
\begin{figure}[htbp]
    \centering
    \includegraphics[width=0.75\textwidth]{figure/img/unexpl_raw_end_both.pdf}
    \caption{Comparison of the Branch-and-Bound tree expansion (unexplored nodes) between the models over time.}
    \label{fig:comp_tree}
\end{figure}
In contrast, the \algoNotation algorithm finishes in an early stage of the branch-and-bound process, with a much smaller tree. The \algoNotation algorithm circumvents the weak LP relaxation by leveraging the Lower Bound Model (LBM) to achieve a much tighter lower bound. The upper bound model finds the integer schedule with an objective value of $-176,843.71$, which perfectly matches the lower bound, the true optimality gap mathematically closes at $0.00\%$.

\section{Analysis of the Optimized Schedule}
The globally optimal solution yields profound operational and commercial insights into scheduling behavior. We systematically extract the activated variables from the exact solution to analyze the underlying operations.

The optimal schedule physically activates exactly $348$ flights, maintaining active service across all $60$ airports in the specified network. The fleet assignment perfectly balances the available sub-types based on demand density: $178$ of the scheduled flights are operated by the higher-capacity A321-271NX, while the remaining $170$ flights are allocated to the A320-214. This heterogeneous deployment ensures that the largest aircraft are strictly assigned to trunk routes with sufficient consolidated demand, minimizing the per-seat-mile cost, while the smaller variants are deployed to efficiently fulfill secondary frequency requirements without excessive capacity spoilage.

Rather than attempting to blindly serve all potential OD pairs, the algorithm exhibits acute profitability-driven selectivity. Out of the $382$ markets, the schedule strategically captures demand from only the $120$ most profitable markets, consciously abandoning structurally unprofitable ones to conserve operating costs. Within these targeted markets, the Sales-Based Linear Program (SBLP) component secures an average market share of $5.42\%$. This is higher than F9's historical market share of $4.85\%$ across these markets, indicating that the optimized schedule is effectively capturing a larger portion of the demand in the most lucrative segments.

Furthermore, the optimization exploits the network's spatial connectivity. The optimal passenger allocation actively utilizes $112$ multi-leg connecting itineraries compared to only $64$ direct itineraries. This heavy reliance on connecting itineraries demonstrates the model's ability to efficiently stitch together fractional demand from disparate markets. By chronologically synchronizing the departure and arrival bounds of physical flights across focus cities, the \algoNotation framework seamlessly captures indirect demand and maximizes the total system-wide operating profit.